Agent memory retrieval requires a selector that delivers relevant items and an accounting of the cost to build them. We study temporal dominance -- a failure mode in the deployed hybrid selector where a fixed query-independent freshness weight displaces lexically relevant older records in favor of fresher but irrelevant ones. We introduce CTLS (Contrastive Tier-Based Lexical Selection), a novel selector that partitions candidates by lexical overlap tier before applying temporal scoring, provably eliminating temporal dominance by construction. On a stratified LongMemEval-S subset of 84 memory episodes under three token budgets, BM25 outperforms the deployed hybrid by -36.1111 percentage points, -29.3651 percentage points, and -33.7302 percentage points respectively (all Holm-corrected). The hybrid\_no\_time ablation recovers most of this gap, confirming the temporal term as the primary cause. CTLS formalizes this fix while preserving temporal decay as a within-tier tiebreaker. No zero-overlap item can outrank a positive-overlap item regardless of freshness gap. CTLS requires no model calls, no embedding model, and no learned parameters; it is a zero-cost selector substitution that leaves construction cost and write amortization unchanged.
